\section{Iterating in 2 subtrees at the same time}

Suppose you are solving a tree problem with dynamic programming. Your DP states look like this: for every vertex $u$ you have a vector $dp[u]$ and the length of this vector is 
at most the number of descendants of $u$. The DP is calculated like this:
\begin{verbatim}
 function calc_dp (u):
     for each child v of u:
         calc_dp(v);
     for each child v of u:
         for each child w of u other than v:
              for i in 0..length(dp[v]):
                  for j in in 0..length(dp[w]):
                      # calculate something 
                      # typically based on dp[v][i] and dp[w][j]
                      # commonly assign to dp[u][i + j]
 
 calc_dp(root);
\end{verbatim}

The complexity of this is $\mathcal{O}(N^2)$ (this is not obvious: when I first saw it I didn't realize it was better than cubic).

Consider the algorithm for a fixed $u$. For each child $v$ of $u$, we associate every position in the array $dp[v]$ with a descendant of $v$. 
The array $dp[v]$ is no longer than the number of descendants of $v$, so we can make sure no two elements in the subtree are associated with the 
same elements in the array. So when you iterate for $i$ and $j$ you are making a pair of 2 vertexs, one in the subtree of $v$ and another in the
subtree of $w$, this pair of vertices won't repeat. So every pair of vertexes in the tree appears only 1 time.

This logic also apply for algorithms like this:
\begin{verbatim}
 function calc_dp (u):
     for each child v of u:
         calc_dp(v);
     tot = 1;
     for each child v of u:
         for i in 0..tot+length(dp[v]):
              tmp[i] = off;
         for i in 0..tot:
              for j in 0..length(dp[v]):
                  tmp[i+j] = calc(dp[u][i],dp[v][j]);
         for i in 0..tot+length(dp[v]):
              dp[u][i] = tmp[i];
         tot += length(dp[v]);
 
 calc_dp(root);
\end{verbatim}
